% capitulos/capitulado-01-fundamentacao.tex
%
% Capítulo 1 do MODELO POR CAPÍTULOS AUTÔNOMOS.
% Traz a fundamentação teórica ampla, o problema de pesquisa, a área
% de estudo, o objetivo geral e — de forma consolidada — todos os
% Objetivos Específicos (OE) e Hipóteses (H) do trabalho. Os demais
% capítulos de conteúdo retomam apenas o OE/H que lhes corresponde,
% com remissão de volta a este capítulo (veja capitulos/capitulado-02-exemplo.tex
% e capitulos/capitulado-03-exemplo.tex).
%
% Todo o conteúdo abaixo é fictício e serve apenas para demonstrar o
% padrão estrutural. Substitua pelo conteúdo real da sua pesquisa.

\section{Estrutura da Tese e Organização dos Capítulos}
\label{sec:estrutura}

Em vez de um formato monolítico tradicional, este trabalho está
estruturado em capítulos autônomos e publicáveis, cada um
correspondente a um Objetivo Específico (OE) e contendo seus próprios
resumo, introdução, objetivo, hipótese (quando aplicável), metodologia
e contribuições esperadas. O \autoref{cap:fundamentacao} apresenta a
fundamentação teórica ampla, o problema de pesquisa, a área de
estudo e — de forma consolidada, na \autoref{sec:oe-h} — todos os
Objetivos Específicos e Hipóteses do trabalho; o último capítulo
integra as contribuições e conclui o trabalho.

Os identificadores de Objetivo Específico (OE) e de Hipótese (H) são
mantidos estáveis e não seguem necessariamente a ordem numérica de
apresentação dos capítulos, evitando ambiguidade em referências
externas já feitas a eles. O texto completo de cada OE e H está
reunido na \autoref{sec:oe-h}; cada capítulo de conteúdo retoma apenas
um resumo de uma linha com remissão de volta a ela.

\begin{table}[H]
  \centering
  \caption{Estrutura da tese e correspondência entre capítulos,
           objetivos específicos e hipóteses (exemplo fictício).}
  \label{tab:estrutura-tese}
  \begin{tabularx}{\textwidth}{cXXc}
    \toprule
    \textbf{Capítulo} & \textbf{Objetivo específico} & \textbf{Foco} & \textbf{Hipótese} \\
    \midrule
    1 & --      & Fundamentação, problema, área de estudo, objetivo geral e \autoref{sec:oe-h} & -- \\
    \thinrule
    2 & OE1     & Título do Capítulo 2 (ver \autoref{cap:capitulo2})   & H1 \\
    \thinrule
    3 & OE2     & Título do Capítulo 3 (ver \autoref{cap:capitulo3})   & -- \\
    \thinrule
    $n$ & --    & Contribuições integradas e conclusão                 & -- \\
    \bottomrule
  \end{tabularx}
  \fonte{Elaborado pelo autor (2025).}
\end{table}

\section{Contextualização}
\label{sec:contextualizacao}

Lorem ipsum dolor sit amet, consectetur adipiscing elit. Pellentesque
habitant morbi tristique senectus et netus et malesuada fames ac turpis
egestas \cite{exemplo:artigo:2024}. O estudo de ecossistemas do Cerrado
brasileiro tem ganhado crescente atenção científica em razão da
singularidade de sua biodiversidade e da acelerada degradação a que está
sujeito \cite{exemplo:livro:2020}.

A contaminação ambiental por metais pesados representa um problema de
saúde pública e ambiental de escala global. Conforme
\citeonline{exemplo:tese:2023}, as concentrações de mercúrio em solos
adjacentes a áreas de garimpo ilegal têm excedido consistentemente os
valores de referência estabelecidos pelos órgãos reguladores
\cite{exemplo:relatorio:2021}.

\section{Problema e Justificativa}
\label{sec:problema}

A fragmentação e a heterogeneidade dos dados de monitoramento
ambiental constituem obstáculo crítico à avaliação da efetividade das
políticas públicas de controle de contaminantes
\cite{exemplo:capitulo:2019}. Estudos anteriores
\cite{exemplo:anais:2023} apontam que a dinâmica hidrológica do
Cerrado cria condicionantes biogeoquímicos ainda pouco compreendidos,
o que reforça a necessidade de um arcabouço integrado — empírico e
computacional — capaz de caracterizar essa dinâmica de forma
sistemática \cite{brasil:lei:2012}.

\section{Objetivo Geral}
\label{sec:objetivo-geral}

Desenvolver e validar um arcabouço integrado para caracterizar
[descreva aqui o objetivo geral do trabalho], articulando um
componente empírico de coleta e análise de dados com um componente
computacional de modelagem e validação.

\section{Área de Estudo}
\label{sec:area-estudo}

A área de estudo compreende [descreva aqui a área de estudo, unidades
de conservação, sítios amostrais ou sistemas comparados]. Quando o
desenho da pesquisa envolver mais de uma unidade ou sítio, descreva
cada um em um parágrafo próprio, explicitando o critério de
comparação adotado (por exemplo, gradiente de exposição, tratamento
\textit{versus} controle, ou série temporal).

\begin{figure}[H]
  \centering
  \fbox{\rule{0pt}{5cm}\rule{10cm}{0pt}}
  \caption{Mapa de localização da área de estudo. Fonte: elaborado
           pelo autor (2025).}
  \label{fig:area-estudo-capitulado}
\end{figure}

\section{Objetivos Específicos e Hipóteses}
\label{sec:oe-h}

Esta seção reúne, num único lugar, todos os Objetivos Específicos
(OE) e Hipóteses (H) do trabalho, na íntegra. Cada capítulo de
conteúdo retoma apenas o OE/H que lhe corresponde, com uma remissão de
volta a esta seção — evitando que o leitor precise localizar cada
objetivo e hipótese espalhados ao longo do texto. A
\autoref{tab:estrutura-tese} resume o mapeamento capítulo–OE–H; esta
seção traz o texto completo.

\subsection{Objetivos Específicos}

\textbf{OE1} (\autoref{cap:capitulo2}) — [Descreva aqui o primeiro
objetivo específico, com a remissão bibliográfica que o fundamenta]
\cite{exemplo:artigo:2024}.

\textbf{OE2} (\autoref{cap:capitulo3}) — [Descreva aqui o segundo
objetivo específico] \cite{exemplo:dissertacao:2022}.

\subsection{Hipóteses}

\textbf{H1} (\autoref{cap:capitulo2}) — [Enuncie aqui a primeira
hipótese testável, com o mecanismo esperado e a referência que a
sustenta] \cite{exemplo:capitulo:2019}.
